\documentclass[11pt,a4paper]{article}
\usepackage[margin=2.4cm]{geometry}
\usepackage[T1]{fontenc}
\usepackage{lmodern}
\usepackage[english]{babel}
\usepackage[hidelinks]{hyperref}
\usepackage{parskip}
\pagestyle{empty}

\begin{document}

Autonomous Teaming Solutions ATS GmbH\\
Taunusstrasse 23\\
80807 Munich\\
Germany

\begin{flushright}
18 August 2026
\end{flushright}

\textbf{Application for Computer Vision Engineer -- Object Detection (m/w/d)}

Dear Autonomous Teaming Hiring Team,

I am applying for the Computer Vision Engineer -- Object Detection position because I am excited by both the technical challenges of autonomous perception and Autonomous Teaming's mission. I believe my experience developing and deploying computer vision systems for difficult real-world data has prepared me to make a strong contribution to your perception team.

My most relevant project was an object-detection model for anatomical structures of the spine used in clinical trials. Because the labeled dataset was small, I designed a custom ViTDet architecture with a DINOv3 backbone using MMDetection and trained it with self-supervised learning on a much larger collection of unlabeled images before supervised fine-tuning. I also designed augmentations to address imbalances in object location, size, and distribution while preserving anatomical constraints. Prior to this work, I worked extensively on a legacy Mask R-CNN detector, gaining practical experience with an established two-stage detection architecture.

I also took ownership of the annotation workflow. I proposed replacing the existing COCO annotation tool with CVAT, labeled data myself, and mentored an intern on using the tooling within his project. This process reinforced how strongly annotation guidelines, label quality, and rapid feedback from model errors influence detection performance.

Beyond medical imaging, I worked as a freelance computer vision engineer on a real-time retail anti-theft system, where I developed depth-map-based camera obstruction detection. This strengthened my experience with real-time inference and event-driven architectures using publish-subscribe messaging. The project, together with my specialized studies, also deepened my understanding of modern detection and tracking methods, including DETR and SORT-based approaches.

I am particularly interested in applying and extending this experience to multi-frame perception and coordinated multi-sensor systems, where motion, occlusion, latency, and limited compute make model and data decisions especially consequential. The mission behind this work is equally important to me. As a Ukrainian citizen, I know personally what defending democracy means, and the opportunity to contribute directly to that cause would mean a great deal to me.

I would welcome the opportunity to discuss how I could contribute to your perception team.

Sincerely,\\
Kostiantyn Pysanyi

\end{document}
